\section{Cài đặt hệ thống}

\subsection{Backend}

Backend được xây dựng bằng FastAPI, đóng vai trò trung tâm trong kiến trúc hệ thống, chịu trách nhiệm
xử lý yêu cầu từ người dùng, quản lý nghiệp vụ, xác thực, phân quyền và kết nối với các dịch vụ như
cơ sở dữ liệu, ứng dụng AI Desktop và hệ thống cảnh báo.

\noindent\textbf{Cách cài đặt:}
\begin{itemize}
\item Khởi tạo môi trường ảo Python bằng lệnh \texttt{python -m venv venv}, sau đó kích hoạt
bằng lệnh \texttt{venv\textbackslash Scripts\textbackslash activate}.
\item Cài đặt các thư viện cần thiết từ file \texttt{requirements.txt} bằng lệnh
\texttt{pip install -r requirements.txt}.
\item Cấu hình các biến môi trường trong file \texttt{.env} dựa theo file mẫu \texttt{.env.example}
để thiết lập kết nối cơ sở dữ liệu, thông tin email và các tham số khác.
\item Khởi động server bằng lệnh \texttt{.\textbackslash run.bat}, Uvicorn sẽ khởi chạy tại
địa chỉ \texttt{localhost} cổng 8000.
\item Kiểm tra các API tại \texttt{http://localhost:8000/docs}.
\end{itemize}

\subsection{Frontend}

Giao diện người dùng gồm ba thành phần: ứng dụng di động Flutter dành cho người dùng cuối,
trang quản trị Nuxt.js dành cho admin và ứng dụng desktop Tkinter tích hợp AI phục vụ giám sát
và nhận diện hành vi.

\textbf{Ứng dụng di động (Flutter):}
\begin{itemize}
\item Cài đặt Flutter SDK và thiết lập môi trường phát triển.
\item Chạy lệnh \texttt{flutter pub get} để cài đặt các gói phụ thuộc.
\item Cấu hình địa chỉ Backend trong file cấu hình của ứng dụng.
\item Khởi chạy ứng dụng trên trình giả lập hoặc thiết bị thực bằng lệnh \texttt{flutter run}.
\end{itemize}

\textbf{Trang quản trị (Nuxt.js):}
\begin{itemize}
\item Cài đặt Node.js và chạy lệnh \texttt{npm install} để cài đặt các gói phụ thuộc.
\item Cấu hình địa chỉ Backend trong file cấu hình của ứng dụng.
\item Khởi chạy ứng dụng bằng lệnh \texttt{npm run dev}, giao diện sẽ chạy tại
\texttt{http://localhost:3000}.
\end{itemize}

\textbf{Ứng dụng desktop AI (Tkinter):}
\begin{itemize}
\item Cài đặt môi trường Python và các thư viện cần thiết từ file
\texttt{requirements.txt}.
\item Đảm bảo mô hình AI đã được huấn luyện và file trọng số
\texttt{.pth} được đặt đúng vị trí trong thư mục dự án.
\item Cấu hình địa chỉ Backend trong file cấu hình của ứng dụng.
\item Khởi chạy ứng dụng bằng lệnh \texttt{.\textbackslash run.bat}.
Khi khởi động, hệ thống sẽ tự động tải mô hình AI từ file
\texttt{.pth} và sẵn sàng thực hiện nhận diện, phân tích dữ liệu
từ camera.
\end{itemize}

\subsection{Cơ sở dữ liệu}

Hệ thống sử dụng PostgreSQL làm cơ sở dữ liệu quan hệ, truy cập thông qua SQLAlchemy kết hợp
AsyncPG theo cơ chế bất đồng bộ. Backend tự động tạo các bảng khi khởi động lần đầu nếu chưa tồn tại.

\noindent\textbf{Cách cài đặt:}
\begin{itemize}
\item Cài đặt PostgreSQL và tạo cơ sở dữ liệu mới cùng tài khoản người dùng với quyền truy cập phù hợp.
\item Chạy các script SQL để tạo bảng và dữ liệu mẫu ban đầu.
\item Cấu hình chuỗi kết nối trong file \texttt{.env}.
\end{itemize}

\subsection{Lưu trữ đám mây}

Hệ thống sử dụng dịch vụ lưu trữ đám mây để quản lý ảnh chụp và video sinh ra trong quá trình giám sát.
Khi ứng dụng AI phát hiện hành vi nguy hiểm, ảnh khung hình được tải lên tự động và đường dẫn được lưu
vào cơ sở dữ liệu. Quyền truy cập tệp được bảo vệ thông qua signed URL có thời hạn.

\noindent\textbf{Cách cài đặt:}
\begin{itemize}
\item Tạo tài khoản và bucket trên dịch vụ lưu trữ đám mây (AWS S3, Google Cloud Storage
hoặc Azure Blob Storage).
\item Cấu hình quyền truy cập và khóa API tương ứng.
\item Cập nhật thông tin kết nối trong file \texttt{.env} và kiểm tra bằng cách tải lên
một tệp thử nghiệm.
\end{itemize}

\subsection{Quy trình khởi động hệ thống}

Các thành phần cần được khởi động theo thứ tự sau:

\begin{itemize}
\item \textbf{Backend}: Kích hoạt môi trường ảo và chạy
\texttt{.\textbackslash run.bat}, đảm bảo cơ sở dữ liệu đã sẵn sàng trước
khi khởi động.

\item \textbf{Frontend}: Khởi chạy ứng dụng di động bằng
\texttt{flutter run} và trang quản trị bằng
\texttt{npm run dev}.

\item \textbf{Ứng dụng desktop AI}: Chạy
\texttt{.\textbackslash run.bat}. Khi khởi động, ứng dụng sẽ tự động
nạp mô hình AI từ file \texttt{.pth}, kết nối với Backend và sẵn sàng
xử lý dữ liệu camera.

\end{itemize}
